% Theorie-Teil der Hausarbeit
% ----------------------------------------------------------------------
% GERUEST-Konventionen (vor Abgabe alle Marker entfernen bzw. aufloesen):
%   SLOT  = hier fehlen eigene 2-4 Saetze zu der im Kommentar beschriebenen
%           Behauptung; Quellenangabe steht schon bereit.
%   NOTIZ = eigene Rohnotiz, noch nicht in Prosa verwandelt.
%   TODO  = Fundstelle (Seitenzahl) noch an der Quelle pruefen.
% Regel: Keine Aussage ohne \parencite, kein \parencite ohne gepruefte
% Fundstelle. Alle SLOT/NOTIZ/TODO muessen vor Abgabe verschwunden sein.

\subsection{Theoretischer Rahmen}

% ----------------------------------------------------------------------
% ABSATZ 1: Begriff und Mechanismus der Konditionalitaet
% (eigene Saetze vorhanden; Seitenzahlen pruefen)
% ----------------------------------------------------------------------
Konditionalität bezeichnet die vom IMF vorgeschriebenen Policy-Reformen \parencite[545]{KentikelenisIMFConditionalityDevelopment2016}.

\textcite{DreherPoliticsIMFConditionality2015} beschreiben Konditionalität als die Notwendigkeit von Regierungen, für den Erhalt notwendiger Kredite von der IMF vorgegebene Policy-Änderungen (Schließen von Haushaltsdefiziten, Erhöhung von Zinsraten, etc.) zu befolgen.

% SLOT 1 — Auszahlungsmechanismus (2-3 eigene Saetze)
% Behauptung: Vergabe und fortlaufende Auszahlung der Kredite sind an die
% Einhaltung der Bedingungen gekoppelt; UEberpruefung erfolgt regelmaessig
% (quartalsweise Reviews); Konditionalitaet ist deshalb kontinuierlich
% verhandelbar, nicht nur eine einmalige Vereinbarung.
% Quelle: \parencite[TODO]{DreherPoliticsIMFConditionality2015}
% NOTIZ (vorhandene Rohnotiz):
%   Die Vergabe und kontinuierliche Auszahlung von IMF-Krediten sind meist
%   daran gekoppelt, ob Regierungen bestimmten Policy-Bedingungen folgen,
%   die jedes Quartal überprüft werden. Diese Arrangements sind bekannt
%   für ihre strenge und kontroverse ökonomische Austerität.

% ----------------------------------------------------------------------
% ABSATZ 2: Politische Oekonomie der Konditionalitaet
% (eigene Saetze vorhanden; NOTIZ-Stueck unten einarbeiten)
% ----------------------------------------------------------------------
Bereits in einer früheren Studie stellten die Autoren fest, dass nicht allen Ländern die gleiche Strenge an Konditionalitäten entgegengebracht wird und der Verdacht bestünde, dass (politische) Nähe zu den Hauptanteilseignern und indirekten Kontrollinstanzen des IMF abgeschwächtere Konditionen verspricht als gegenüber Ländern mit weniger Relevanz \parencite{DesaiGlobalGovernanceMultipolar2011}.

% SLOT 2 — Zwei Wege politischer Einflussnahme (2-3 eigene Saetze)
% Behauptung: Der Literaturstrang um Stone unterscheidet zwei Mechanismen:
% (a) mildere Sanktionen bei Nicht-Einhaltung fuer strategisch wichtige
% Laender, (b) geringere Zahl und Schwere der Bedingungen selbst fuer
% Laender in der Gunst der Hauptanteilseigner (USA, Japan, Deutschland,
% Frankreich, Grossbritannien).
% Quelle: \parencite[TODO]{DreherPoliticsIMFConditionality2015} und dortige
% Verweise auf Stone 2002/2004/2008.
% ACHTUNG: Stone 2002/2004/2008 sind NICHT in literatur.bib — entweder
% Beitraege anlegen und selbst lesen oder die Passage als Verweis ueber
% \parencite{DreherPoliticsIMFConditionality2015} formulieren.
% NOTIZ (vorhandene Rohnotiz):
%   1. Möglichkeit: Konditionalität ist weniger strikt, wenn Kredite an
%   Ländern vergeben werden, die strategisch wichtig für die
%   Hauptanteilseigner des IMF sind (siehe z.B. Stone 2002 und 2004:
%   Strafen bei Nicht-Einhaltung fallen signifikant schwächer aus für
%   Länder, die von Bedeutung für die USA sind).
%   2. Möglichkeit: Die Konditionen selbst fallen in Anzahl und Schwere
%   geringer aus für Länder, die in der Gunst der mächtigsten Länder des
%   IMF stehen (siehe Stone 2008). Dreher et al. argumentieren, dass
%   diese Möglichkeiten erklären würden, warum Länder, die zugunsten der
%   USA, Japan, Deutschlands, Frankreichs und Großbritanniens stimmen,
%   eine größere Chance haben, IMF-Kredite zu erhalten.

% ----------------------------------------------------------------------
% ABSATZ 3: UNSC-Mitgliedschaft als quasi-exogenes Messinstrument
% (eigene Saetze vorhanden)
% ----------------------------------------------------------------------
Es bedarf daher eines Messinstruments der (politischen und ökonomischen) Relevanz, das zu keinem Endogenitätsproblem führt: die UNSC-Mitgliedschaft als quasi-exogenes Messinstrument.

Dreher et al. schlagen die Mitgliedschaft im UN-Sicherheitsrat vor, da temporäre Mitglieder größere Finanzierung von unterschiedlichen Quellen, darunter auch dem IMF, erhalten. Dieses Geld wird vermutlich dafür verwendet, Wahlzugeständnisse zu erkaufen, z. B. in Bezug auf Voten, Sanktionen und militärische Interventionen. Die Hauptanteilseigner des IMF (USA, Japan, Deutschland, Frankreich und Großbritannien) möchten ihren Einfluss auf den Sicherheitsrat erhöhen und bündeln die Kosten dieser Einflussnahme durch Kreditvergabe des IMF. Indem Länder des Globalen Südens während ihrer Mitgliedschaft Wahlzugeständnisse garantieren, erhalten sie im Gegenzug bessere Bedingungen bei der Kreditvergabe \parencite[121]{DreherPoliticsIMFConditionality2015}.

% ----------------------------------------------------------------------
% ABSATZ 4: Abhaengigkeitstheoretische Einordnung
% (noch leer — hier eigene Saetze schreiben; Quellen bereits gelesen)
% ----------------------------------------------------------------------
% SLOT 3 — Grundlogik der Abhaengigkeitstheorie (2-3 eigene Saetze)
% Behauptung: Kernannahme der Dependency Theory zum Verhaeltnis von
% Zentrum und Peripherie in der Weltwirtschaft (ungleiche
% Austauschbeziehungen, strukturelle Abhaengigkeit der Peripherie von
% Kapital und Krediten des Zentrums).
% Quellen: \parencite[TODO]{AntunesDeOliveiraDependencyTheory2024}
% (Kernquelle, gelesen), \parencite[TODO]{MuchhalaExaminingFeaturesEconomic2025}
% (gelesen; oekonomische Kolonialitaet, epistemische Singularitaet).

% SLOT 4 — Anschluss an IMF-Konditionalitaet (2-3 eigene Saetze)
% Behauptung: Aus abhaengigkeitstheoretischer Sicht ist IMF-Konditionalitaet
% ein Instrument, mit dem Zentren die Politikoekonomie der Peripherie
% formen; ungleiche Strenge zwischen Laendern ist dann kein Anomalie,
% sondern Ausdruck asymmetrischer Machtbeziehungen.
% Quellen: \parencite[TODO]{MuchhalaExaminingFeaturesEconomic2025},
% \parencite{deOliveira}.

% SLOT 5 — Bruecke zur eigenen Fragestellung (2-3 eigene Saetze)
% Behauptung: Wenn politische Relevanz (UNSC) und oekonomische
% Machtressourcen (Rohstoffabhaengigkeit) die Strenge von Bedingungen
% moderieren, verbindet die Arbeit die politoekonomische Empirie von
% Dreher et al. mit der strukturellen Perspektive der Dependency Theory;
% daraus folgt die Frage, ob rohstoffreiche Laender aehnliche
% Entgegenkommen erfahren wie UNSC-Mitglieder (Anschluss an H2).

\subsection{Hypothesen}

% SLOT 6 — H1 in eigenen Worten (1-2 Saetze)
% Behauptung: Temporaere UNSC-Mitglieder erhalten waehrend ihrer
% Mitgliedschaft weniger Bedingungen pro Quartal als vergleichbare
% Nicht-Mitglieder. Richtung der Erwartung begruenden (Nettoeffekt unter
% Kontrolle von Programmdauer und Wirtschaftsstruktur; vgl. Absatz zur
% deskriptiven Evidenz in daten-methodik.tex).
% Modell im Datensatz (Phase B): avgcondtype_count ~ unsc3 + Kontrollen.
% Messfenster von unsc3: Ein Land-Jahr t zaehlt als UNSC-"Mitglied",
% wenn das Land in t ODER t+1 temporAer Mitglied ist (DSV-Regel, Wahljahr
% inklusive) — also Wahljahr + 2 Amtsjahre, nicht nur die Amtsjahre.
% Belege: CODEBOOK_MONA.md Abschnitt 10.5 (dort korrigiert) und
% code/shared/data_prep/build_unsc_dsv_rule.R; die Operationalisierung
% in daten-methodik.tex bei den unabhaengigen Variablen erwaehnen.
% Quelle: \parencite[TODO]{DreherPoliticsIMFConditionality2015}
% NOTIZ (vorhandene Rohnotiz — Achtung: enthaelt ERGEBNISSE, die in
% ergebnisse.tex gehoeren; hier nur die inhaltliche Erwartung behalten):
%   Die Originalstudie stellt die Hypothese auf, dass temporäre
%   UNSC-Mitglieder weniger Konditionen erhalten als andere Länder, die
%   ebenfalls an IMF-Programmen teilnehmen. Ergebnisse: 2,5 weniger
%   Konditionen. Ein Programm hat durchschnittlich 80 Prozent
%   Policy-Konditionen, macht das ca. 30 Prozent aus. Vor allem betroffen
%   sind „Prio Actions“ (um Kredite zu erhalten) und „Performance
%   Criteria“ (um Kreditauszahlungen fortzuführen). Schwächere Ergebnisse
%   für „Structural Benchmarks“ und allgemeinen Konditionsumfang. Die
%   Autoren merken an, dass der Umfang an Policy-Bereichen, die von
%   „Performance Criteria“ abgedeckt sind, an sich bereits enger ist für
%   UNSC-Mitglieder.

% SLOT 7 — H2 in eigenen Worten (1-2 Saetze)
% Behauptung: Die theoretically erwartete Moderation durch
% Rohstoffabhaengigkeit formulieren (Richtung begruenden mit SLOT 4/5).
% Modell im Datensatz (Phase B): avgcondtype_count ~ unsc3 * resource_dep + Kontrollen.
% (unsc3 mit demselben 3-Jahres-Fenster wie in SLOT 6: Wahljahr + 2
% Amtsjahre, DSV-Regel t | t+1.)
% Quellen: eigene Argumentation aus SLOT 4/5, gestuetzt auf
% \parencite[TODO]{DreherPoliticsIMFConditionality2015} und
% \parencite[TODO]{MuchhalaExaminingFeaturesEconomic2025}.

% SLOT 8 — Abgrenzung H3/H4 (1-2 eigene Saetze, optional)
% Behauptung: Kurz erwaehnen, dass die Laenderprofile (H3) und die
% Kompositionsanalyse (H4) explorativ bleiben (Begruendung ausfuehrlich
% in daten-methodik.tex, dortiger Absatz zu Arrangement-Typen und
% Policy-Areas).
